\begin{table}[htbp]
\centering
\small
\caption{Bloco 2, Problema E3 (openai/gpt-4o-mini): aprendizado em contexto por estratégia de seleção e n.}
\label{tab:bloco2_E3_faseE_estrategia_x_nshot_gpt4mini}
\begin{tabular}{lrrrrrrrr}
\toprule
 & \multicolumn{2}{c}{Fácil} & \multicolumn{2}{c}{Difícil} & \multicolumn{2}{c}{Mista} & \multicolumn{2}{c}{Aleatória} \\ \cmidrule(lr){2-3} \cmidrule(lr){4-5} \cmidrule(lr){6-7} \cmidrule(lr){8-9}
n & Acur.\ vs espec. & $\kappa$ & Acur.\ vs espec. & $\kappa$ & Acur.\ vs espec. & $\kappa$ & Acur.\ vs espec. & $\kappa$ \\
\midrule
zero-shot (único valor) & 53,6 (50,7–55,3) & 0,096 &  &  &  &  &  &  \\
4-shot & 76,0 (70,5–79,5) & 0,529 & 53,2 (50,0–56,2) & 0,053 & 88,8 (80,1–93,2) & 0,777 & 86,1 (71,9–93,8) & 0,724 \\
10-shot & 75,5 (70,0–82,1) & 0,519 & 56,9 (51,4–63,6) & 0,152 & 84,8 (78,6–87,9) & 0,695 & 91,7 (88,6–95,0) & 0,833 \\
20-shot & 80,8 (78,5–83,1) & 0,620 & 75,1 (66,9–90,8) & 0,509 & 90,0 (86,2–93,1) & 0,800 & 93,9 (88,5–99,2) & 0,879 \\
40-shot & 83,3 (79,1–87,3) & 0,669 & 97,6 (97,3–98,2) & 0,951 & 93,6 (88,2–97,3) & 0,872 & 94,7 (89,1–99,1) & 0,895 \\
\bottomrule
\multicolumn{9}{p{0.95\linewidth}}{\footnotesize Especialista E3, W = [1,0; 1,0]: rotula o Problema E pelo centróide mais próximo; o LLM recebe n exemplos escolhidos pela estratégia e classifica os demais. Acur.\ vs espec.\ = acurácia (LLM vs especialista): fração dos pontos de teste em que o LLM reproduz o rótulo do especialista. Valores em \% (mín–máx); $\kappa$ = Cohen's $\kappa$ (faixas de Landis e Koch, 1977). Médias das sementes e sorteios. Em zero-shot não há exemplo a selecionar: valor único.} \\
\end{tabular}
\end{table}
